% Chapter 6 -- Determinants
\chapter{Determinants}\label{ch:determinants}

Every question this book has asked about a square matrix --- does
$Ax=\mathbf b$ have a unique solution? is $A$ invertible? does the map
collapse space? --- turns out to be one number's job. The determinant is
that number: a single scalar distilled from $n^2$ entries that answers all
of them at once, and geometrically measures how much a linear map stretches
volume. After this chapter you will be able to compute determinants three
ways, know which way to reach for in a computation, and read $\det A$ as a
signed volume rather than an abstract formula. One honest warning before we
start: determinants are \emph{theoretical} heavyweights and
\emph{computational} lightweights --- you almost never need one in code,
and this chapter explains why.

\section{Three definitions that agree}\label{sec:det-defs}

\begin{definition}[Determinant by cofactors]\index{determinant}\index{cofactor expansion}
For $A\in\mathbb R^{n\times n}$, the \emph{$(i,j)$-minor} $M_{ij}$ is the
determinant of the matrix left by deleting row $i$ and column $j$; the
\emph{cofactor} is $C_{ij}=(-1)^{i+j}M_{ij}$. Then
\[
\det A=\sum_{j=1}^{n}a_{1j}\,C_{1j}
\qquad\text{(expansion along row 1)},
\]
with $\det[a]=a$ starting the recursion. The same sum along \emph{any} row
or column gives the same value.
\end{definition}

\begin{definition}[Determinant by pivots]\index{pivot determinant}
If elimination reduces $A$ to upper-triangular $U$ using only row-add
operations and $s$ row swaps, then
\[
\det A=(-1)^{s}\,u_{11}u_{22}\cdots u_{nn},
\]
the product of the pivots with one minus sign per swap. If no nonzero pivot
exists in some column, $\det A=0$.
\end{definition}

\begin{definition}[Determinant as signed volume]\index{signed volume}
$\det A$ is the signed volume of the parallelepiped spanned by the columns
of $A$ --- equivalently, the factor by which the linear map $A$ scales
$n$-dimensional volume, negative when orientation flips.
\end{definition}

That three definitions give one function is the content of the chapter;
Figure~\ref{fig:det-scaling} shows the volume meaning directly.

\begin{example}[One 3$\times$3, two ways]\label{ex:det3}
\textbf{Strategy.} Cofactors along a row with zeros keep the expansion
short; elimination needs no swaps here, so the pivot product is a pure
product. Compare.

\textbf{Solution.} Let
$A=\begin{bmatrix}2&1&0\\1&3&1\\0&1&2\end{bmatrix}$.
Cofactors along row 1 (the zero in position $(1,3)$ removes a term):
\[
\det A
=2\begin{vmatrix}3&1\\1&2\end{vmatrix}
-1\begin{vmatrix}1&1\\0&2\end{vmatrix}
+0
=2(6-1)-1(2-0)=10-2=8 .
\]
By elimination: $R_2\leftarrow R_2-\tfrac12R_1$, then
$R_3\leftarrow R_3-\tfrac{2}{5}R_2$:
\[
\begin{bmatrix}2&1&0\\0&5/2&1\\0&0&8/5\end{bmatrix},
\qquad \det A=2\cdot\tfrac52\cdot\tfrac85=8 .
\]
Same answer, $8$ --- and the second route produced it with three
multiplies instead of a tree of $2\times2$ determinants. That asymmetry
only gets worse with size (checkpoint below).
\end{example}

The pivot-product route is an algorithm, not just a trick --- here it is in
code, swap counting included:

\begin{lstlisting}[caption={\ttfamily code/ch06\_pivot\_det.py --- $\det$ as
signed pivot product, verified against \ttfamily np.linalg.det on the
Example~\ref{ex:det3} matrix and on a matrix that forces a row swap.}]
import numpy as np
np.set_printoptions(precision=4, suppress=True)

def pivot_det(A):
    """det via elimination: product of pivots times (-1)^swaps."""
    A = A.astype(float).copy()
    n = A.shape[0]
    swaps = 0
    det = 1.0
    for j in range(n):
        piv = np.argmax(np.abs(A[j:, j])) + j
        if abs(A[piv, j]) < 1e-12:
            return 0.0
        if piv != j:
            A[[j, piv]] = A[[piv, j]]
            swaps += 1
        det *= A[j, j]
        A[j+1:, j:] -= np.outer(A[j+1:, j] / A[j, j], A[j, j:])
    return det * (-1.0) ** swaps

# 3x3 from the chapter: det = 8 by cofactors AND by pivots
A = np.array([[2., 1., 0.],
              [1., 3., 1.],
              [0., 1., 2.]])
print("matrix:")
print(A)
print("pivot_det =", pivot_det(A))
print("np.linalg.det =", np.linalg.det(A))

# swap-count check on a matrix needing a swap
B = np.array([[0., 2., 1.],
              [1., 0., 1.],
              [2., 1., 0.]])
print("swap-needed matrix pivot_det =", pivot_det(B),
      "np.linalg.det =", np.linalg.det(B))
\end{lstlisting}

\begin{lstlisting}[caption={Output of the pivot determinant. Notice the
swap-counted version matches the library on both --- including the matrix
whose zero top-left forces a swap ($\det=5$).}]
matrix:
[[2. 1. 0.]
 [1. 3. 1.]
 [0. 1. 2.]]
pivot_det = 8.0
np.linalg.det = 8.000000000000002
swap-needed matrix pivot_det = 5.0 np.linalg.det = 4.999999999999999
\end{lstlisting}

\begin{exercise}[easy]\label{ex:det-e1}
Compute $\det\begin{bmatrix}3&-1\\2&4\end{bmatrix}$ by the $2\times2$
formula.
\begin{solution}
$3\cdot4-(-1)\cdot2=12+2=14$.
\end{solution}
\end{exercise}

\section{Properties: product, transpose, row operations}\label{sec:det-props}

\begin{theorem}[Working properties]\label{th:detprops}
For $A,B\in\mathbb R^{n\times n}$:
\begin{itemize}
\item $\det(A^{T})=\det A$.
\item $\det(AB)=\det A\,\det B$ --- and $\det(A^{-1})=1/\det A$ when $A$ is
invertible.
\item Row-add ($R_i\leftarrow R_i+cR_j$) leaves $\det$ unchanged; a row
swap flips the sign; scaling one row by $c$ scales $\det$ by $c$, so
$\det(cA)=c^{n}\det A$ (not $c\det A$).
\end{itemize}
\end{theorem}

\begin{example}[The product rule, numerically and in outline]\label{ex:detprod}
\textbf{Strategy.} Check on a pair of $2\times2$s first --- catching a sign
or arithmetic slip is easy there --- then see why the property is
structural rather than coincidental.

\textbf{Solution.} Let $A=\begin{bmatrix}2&1\\0&3\end{bmatrix}$ and
$B=\begin{bmatrix}1&2\\1&1\end{bmatrix}$. Then $\det A=6$, $\det B=-1$,
and
\[
AB=\begin{bmatrix}3&5\\3&3\end{bmatrix},
\qquad \det(AB)=9-15=-6=6\cdot(-1) .
\]
Why it holds (sketch): write $A$ and $B$ as chains of elementary row
operations (ch.~2 --- every invertible matrix is such a product). Each
elementary factor multiplies a determinant by a known factor --- $c$ for a
scaling, $-1$ for a swap, $1$ for a row-add --- and the factors multiply
through the chain. Volume reading: applying $B$ then $A$ scales volume
first by $\det B$, then by $\det A$; the composition scales by their
product.
\end{example}

\begin{checkpoint}
Which is bigger work for a $20\times20$: cofactor expansion or pivot
product? Cofactors cost $\sim n!$ operations; elimination $\sim n^3/3$.
($20!\approx2.4\times10^{18}$ vs $\approx2700$ --- this is why
``determinant'' in software is a pivot product.)
\end{checkpoint}

\section{Determinant as signed volume}\label{sec:det-volume}

A linear map carries the unit square (unit cube) to a parallelogram
(parallelepiped) whose edges are the columns of $A$. The signed area of
that image is exactly the determinant --- the sign recording whether the
map preserved or flipped orientation
(Figures~\ref{fig:det-scaling}--\ref{fig:det-collapse}).

\begin{figure}[htbp]
\centering
\begin{tikzpicture}[scale=0.75, >=stealth]
  % left: unit square grid
  \begin{scope}
    \draw[step=1.0,graybg!60] (0,0) grid (4,4);
    \draw[->,gray] (0,0)--(4.5,0); \draw[->,gray] (0,0)--(0,4.5);
    \fill[accentsoft!40] (0,0)--(1,0)--(1,1)--(0,1)--cycle;
    \draw[accent,thick] (0,0)--(1,0)--(1,1)--(0,1)--cycle;
    \node at (2,-0.7) {\small before: area $=1$};
  \end{scope}
  \draw[->,very thick,ink] (5,2)--(6.4,2) node[midway,above]{$A$};
  % right: mapped parallelogram, A=[[2,1],[1,1.5]]
  \begin{scope}[xshift=7.5cm]
    \draw[step=1.0,graybg!60] (0,0) grid (4,4);
    \draw[->,gray] (0,0)--(4.5,0); \draw[->,gray] (0,0)--(0,4.5);
    \fill[accent!35] (0,0)--(2,1)--(3,2.5)--(1,1.5)--cycle;
    \draw[accent,thick] (0,0)--(2,1)--(3,2.5)--(1,1.5)--cycle;
    \draw[->,good,thick] (0,0)--(2,1) node[midway,below left=-1pt]{\tiny $A\mathbf e_1$};
    \draw[->,amber,thick] (0,0)--(1,1.5) node[midway,above left=-1pt]{\tiny $A\mathbf e_2$};
    \node at (2,-0.7) {\small after: area $=|\det A|=2$};
  \end{scope}
\end{tikzpicture}

\caption{The map $A=\begin{bmatrix}2&1\\1&1.5\end{bmatrix}$ sends the unit
square to a parallelogram of area $|\det A|=2$. Notice the edges of the
image are the \emph{columns} of $A$. Regenerate:
\ttfamily figures/ch06\_volume\_scaling.tex.}
\label{fig:det-scaling}
\end{figure}

\begin{example}[Volume of a parallelepiped]\label{ex:detvol}
\textbf{Strategy.} Three edges from a common vertex make the columns;
the signed volume is the determinant, and its sign tells whether the
three edges form a right- or left-handed triple.

\textbf{Solution.} Edges $\mathbf u=(1,0,0)$, $\mathbf v=(1,1,0)$,
$\mathbf w=(0,1,2)$:
\[
\det\begin{bmatrix}1&1&0\\0&1&1\\0&0&2\end{bmatrix}=1\cdot1\cdot2=2,
\]
triangular, so the product of the diagonal. Volume $=2$, orientation
preserved (positive). Reversing the order of $\mathbf u$ and $\mathbf v$
(column swap) gives $-2$: same shape, opposite handedness.
\end{example}

\begin{figure}[htbp]
\centering
\begin{tikzpicture}[scale=0.8, >=stealth]
  \begin{scope}
    \draw[->,gray] (-0.3,0)--(2.6,0); \draw[->,gray] (0,-0.3)--(0,2.6);
    \draw[->,accent,very thick] (0,0)--(1.9,0.4) node[right]{\tiny $\mathbf u$};
    \draw[->,accent,very thick] (0,0)--(0.7,1.8) node[above]{\tiny $\mathbf v$};
    \draw[->,gray] (1.6,0.45) arc (14:65:0.9);
    \node at (1.7,1.05) {\tiny ccw};
    \node at (1.3,-0.7) {\small $\det>0$: ccw order kept};
  \end{scope}
  \begin{scope}[xshift=5.5cm]
    \draw[->,gray] (-0.3,0)--(2.6,0); \draw[->,gray] (0,-0.3)--(0,2.6);
    \draw[->,amber,very thick] (0,0)--(1.9,0.4) node[right]{\tiny $\mathbf u$};
    \draw[->,amber,very thick] (0,0)--(0.7,-1.8) node[below]{\tiny $F\mathbf v$};
    \draw[->,gray] (1.6,-0.45) arc (-14:-65:0.9);
    \node at (1.7,-1.05) {\tiny cw};
    \node at (1.3,-2.4) {\small $\det<0$: orientation flipped};
  \end{scope}
\end{tikzpicture}

\caption{Orientation and the sign of $\det$. Notice a positive determinant
keeps the counterclockwise ordering of the image edges; a reflection
reverses it and $\det$ goes negative. Regenerate:
\ttfamily figures/ch06\_orientation.tex.}
\label{fig:det-orientation}
\end{figure}

The code tests the scaling claim on shapes that are \emph{not} the unit
square --- the factor should be the same for any region.

\begin{lstlisting}[caption={\ttfamily code/ch06\_volume\_scaling.py ---
areas of three different shapes before/after a $2\times2$ map, then a
parallelepiped through a $3\times3$.}]
import numpy as np
np.set_printoptions(precision=4, suppress=True)
np.random.seed(7)

def poly_area(P):
    x, y = P[:, 0], P[:, 1]
    return 0.5 * abs(np.sum(x * np.roll(y, -1) - y * np.roll(x, -1)))

A = np.array([[2., 1.], [1., 1.5]])
shapes = {
    "unit square": np.array([[0, 0], [1, 0], [1, 1], [0, 1.]]),
    "rectangle": np.array([[0, 0], [2, 0], [2, 0.5], [0, 0.5]]),
    "triangle": np.array([[0, 0], [1, 0], [0.5, 1.2]]),
}
print("A =")
print(A)
print("|det A| =", abs(np.linalg.det(A)))
for name, S in shapes.items():
    before = poly_area(S)
    after = poly_area(S @ A.T)
    print(f"{name:12s} area {before:.4f} -> {after:.4f}  "
          f"ratio {after/before:.4f}")

B = np.random.randn(3, 3)
E = np.array([[1., 0.5, 0.], [0., 1., 0.5], [0.5, 0., 1.]])
v0 = abs(np.linalg.det(E))
v1 = abs(np.linalg.det(B @ E))
print("3D parallelepiped: |det E| =", v0,
      " |det(B E)| =", v1, " ratio =", v1 / v0,
      " |det B| =", abs(np.linalg.det(B)))
\end{lstlisting}

\begin{lstlisting}[caption={Output of the scaling experiment. Notice every
shape scales by exactly $|\det A|=2$ --- the factor is a property of the
map, not of the region. In 3D the ratio equals $|\det B|$ to the last
digit shown, which is also the product rule at work.}]
A =
[[2.  1. ]
 [1.  1.5]]
|det A| = 2.0
unit square  area 1.0000 -> 2.0000  ratio 2.0000
rectangle    area 1.0000 -> 2.0000  ratio 2.0000
triangle     area 0.6000 -> 1.2000  ratio 2.0000
3D parallelepiped: |det E| = 1.125   |det(B E)| = 1.3290513746968733   ratio = 1.181378999730554   |det B| = 1.181378999730554
\end{lstlisting}

\section{\texorpdfstring{$\det A=0$ means collapse}{det A = 0 means collapse}}\label{sec:det-zero}

If the parallelepiped collapses flat --- the columns are dependent --- the
volume is zero. So $\det A=0$, dependent columns, non-invertible, singular:
four ways to say the same thing (the equivalence built through
chapters 1--5). Figure~\ref{fig:det-collapse} is the picture.

\begin{figure}[htbp]
\centering
\begin{tikzpicture}[scale=0.75, >=stealth]
  \begin{scope}
    \draw[step=1.0,graybg!60] (0,0) grid (4,4);
    \draw[->,gray] (0,0)--(4.5,0); \draw[->,gray] (0,0)--(0,4.5);
    \fill[accentsoft!40] (0,0)--(1,0)--(1,1)--(0,1)--cycle;
    \draw[accent,thick] (0,0)--(1,0)--(1,1)--(0,1)--cycle;
    \node at (2,-0.7) {\small volume $1$};
  \end{scope}
  \draw[->,very thick,ink] (5,2)--(6.4,2) node[midway,above]{$A$};
  \begin{scope}[xshift=7.5cm]
    \draw[step=1.0,graybg!60] (0,0) grid (4,4);
    \draw[->,gray] (0,0)--(4.5,0); \draw[->,gray] (0,0)--(0,4.5);
    \draw[accent,very thick] (0,0)--(3.5,3.5);
    \fill[accent] (1,1) circle (2pt) (2.5,2.5) circle (2pt);
    \node[above left=-2pt] at (1,1) {\tiny $A\mathbf e_1$};
    \node[above left=-2pt] at (2.5,2.5) {\tiny $A\mathbf e_2$};
    \node at (2,-0.7) {\small square $\to$ segment: $\det=0$};
  \end{scope}
\end{tikzpicture}

\caption{$\det A=0$ is geometric collapse: the square's image is a segment.
Notice both column images land on one line --- dependent columns are the
\emph{cause}, zero volume is the \emph{symptom}. Regenerate:
\ttfamily figures/ch06\_collapse.tex.}
\label{fig:det-collapse}
\end{figure}

The trap is the word \emph{near}. A small determinant is not a small
collapse --- it is not a collapse at all:

\begin{lstlisting}[caption={\ttfamily code/ch06\_near\_singular.py ---
four $2\times2$s from ``det $=1$'' to ``det $=0$'', with the condition
number alongside. Watch which quantity actually tracks invertibility.}]
import numpy as np
np.set_printoptions(precision=6, suppress=True)

A = np.array([[1., 0.], [0., 1.]])          # det = 1
B = np.array([[100., 0.], [0., 0.001]])     # det = 0.1 -- small!
C = np.array([[1., 1.], [1., 1.0001]])      # det = 1e-4 -- nearly singular
D = np.array([[1., 1.], [1., 1.]])          # det = 0 -- singular
for name, M in [("I", A), ("diag(100,0.001)", B),
                ("near-dep", C), ("dup-rows", D)]:
    d = np.linalg.det(M)
    c = np.linalg.cond(M)
    inv_ok = "yes"
    try:
        _ = np.linalg.inv(M)
    except np.linalg.LinAlgError:
        inv_ok = "NO"
    print(f"{name:16s} det={d:12.6f}  cond={c:12.4f}  invertible={inv_ok}")

B2 = B / np.abs(B).max()
print("rescaled diag(1, 1e-5): det =", np.linalg.det(B2),
      " cond =", np.linalg.cond(B2))
\end{lstlisting}

\begin{lstlisting}[caption={Output of the near-singular test. Notice
``det $=0.1$'' does not mean safe --- the matrix with det $0.1$ has
condition $10^{5}$, and only the truly singular matrix fails inversion.
Determinant answers \emph{is} it invertible; condition number answers
\emph{how close}.}]
I                det=    1.000000  cond=      1.0000  invertible=yes
diag(100,0.001)  det=    0.100000  cond= 100000.0000  invertible=yes
near-dep         det=    0.000100  cond=  40002.0001  invertible=yes
dup-rows         det=    0.000000  cond=139378947902558128.0000  invertible=NO
rescaled diag(1, 1e-5): det = 9.999999999999997e-06  cond = 99999.99999999999
\end{lstlisting}

\begin{exercise}[medium]\label{ex:det-m1}
$\det A=6$ and $\det B=-2$ for $3\times3$ matrices. Compute
$\det(ABA^{-1})$, $\det(A^{T}B)$, and $\det(2A)$.
\begin{solution}
$\det(ABA^{-1})=\det B=-2$ (similarity). $\det(A^{T}B)=\det A\det B=-12$.
$\det(2A)=2^{3}\cdot6=48$ --- the exponent is $n$, not $1$.
\end{solution}
\end{exercise}

\begin{exercise}[hard]\label{ex:det-h1}
Prove the pivot-product formula for a $2\times2$: eliminating
$R_2\leftarrow R_2-(c/a)R_1$ in $\begin{bmatrix}a&b\\c&d\end{bmatrix}$
($a\ne0$) leaves determinant $a(d-bc/a)=ad-bc$; then show swapping the rows
of $\begin{bmatrix}a&b\\c&d\end{bmatrix}$ flips the sign.
\begin{solution}
After $R_2\leftarrow R_2-(c/a)R_1$ the matrix is
$\begin{bmatrix}a&b\\0&d-bc/a\end{bmatrix}$; row-add leaves det unchanged,
triangular det is the diagonal product $a(d-bc/a)=ad-bc$. Swapped matrix
$\begin{bmatrix}c&d\\a&b\end{bmatrix}$ has det $cb-da=-(ad-bc)$.
\end{solution}
\end{exercise}

\section{Cramer's rule, a curiosity}\label{sec:det-cramer}

\begin{theorem}[Cramer's rule]\index{Cramer's rule}
If $A$ is invertible, the unique solution of $Ax=\mathbf b$ is
\[
x_j=\frac{\det A_j}{\det A},
\qquad
A_j=\text{$A$ with column $j$ replaced by $\mathbf b$}.
\]
\end{theorem}

Elegant, correct, and useless for computation: it computes $n+1$
determinants where elimination computes the answer directly. Its real
life is theoretical --- it proves the solution is a rational function of
the entries, and it hides inside the formula
$A^{-1}=\operatorname{adj}(A)/\det A$. Meet it so you can say you have met
it; then go back to elimination.

\begin{figure}[htbp]
\centering
\begin{tikzpicture}[scale=0.95]
  \node[minimum size=0.95cm,draw=gray,fill=accentsoft!25] at (0,0) {$+$};
  \node[minimum size=0.95cm,draw=gray,fill=amberbg!40]     at (1,0) {$-$};
  \node[minimum size=0.95cm,draw=gray,fill=accentsoft!25] at (2,0) {$+$};
  \node[minimum size=0.95cm,draw=gray,fill=amberbg!40]     at (0,-1) {$-$};
  \node[minimum size=0.95cm,draw=gray,fill=accentsoft!25] at (1,-1) {$+$};
  \node[minimum size=0.95cm,draw=gray,fill=amberbg!40]     at (2,-1) {$-$};
  \node[minimum size=0.95cm,draw=gray,fill=accentsoft!25] at (0,-2) {$+$};
  \node[minimum size=0.95cm,draw=gray,fill=amberbg!40]     at (1,-2) {$-$};
  \node[minimum size=0.95cm,draw=gray,fill=accentsoft!25] at (2,-2) {$+$};
  \node[left] at (-0.6,-1) {\small sign of position $(i,j)$: $(-1)^{i+j}$};
\end{tikzpicture}

\caption{The cofactor sign pattern $(-1)^{i+j}$ for $3\times3$. Notice it
is a checkerboard --- diagonally adjacent positions share a sign, which is
why you can expand along any row. Regenerate:
\ttfamily figures/ch06\_checkerboard.tex.}
\label{fig:det-checker}
\end{figure}

\section{Hand derivation}\label{sec:handderiv6}

Derive the $2\times2$ formula from elimination, then watch a row swap flip
the sign. For $A=\begin{bmatrix}a&b\\c&d\end{bmatrix}$ with $a\ne0$:
$R_2\leftarrow R_2-(c/a)R_1$ gives
$\begin{bmatrix}a&b\\0&d-bc/a\end{bmatrix}$ --- upper triangular,
determinant the pivot product $a\cdot(d-bc/a)=ad-bc$. If $a=0$ the first
step swaps the rows; the swap contributes a minus sign and the formula
still lands on $ad-bc$ (check the boundary case $a=0$, $c\ne0$
directly: $ad-bc=-bc$, and the swap-version gives
$cb-da=-(ad-bc)=bc-ad$ --- consistent once the sign travels with it).
Complete the check every derivation owes: the formula is antisymmetric in
the two rows ($cb-da=-(ad-bc)$) and vanishes on dependent rows
($a/c=b/d$ gives $ad-bc=0$).

\section{Common mistakes}

\begin{tripwire}[{Adding determinants: $\det(A+B)\ne\det A+\det B$}]
Determinants multiply; they do not add. Counterexample in one line:
$\det I+\det I=2$ but $\det(2I)=2^{n}$. If you ever split a determinant
over a sum, you have produced a wrong number.
\end{tripwire}

\begin{tripwire}[{Sign bookkeeping on swaps}]
One row swap flips the sign; two swaps flip it back. Learners reliably
lose this --- when you count swaps, count them on paper, and write the
$(-1)^{s}$ factor into the formula before the pivot product, not after
``remembering at the end.''
\end{tripwire}

\begin{tripwire}[{Computing with Cramer's rule}]
Cramer is $n+1$ determinants per solve; elimination is one factorization.
In code, Cramer's rule on a $10\times10$ is already thousands of times
slower, and less stable. The formula is for theorems, not for numbers.
\end{tripwire}

\begin{tripwire}[{Reading ``small determinant'' as ``singular''}]
$\det A=10^{-6}$ is small but nonzero: $A$ is invertible, full stop.
``Nearly singular'' is a statement about \emph{conditioning}
(Section~\ref{sec:det-zero}), not about the determinant crossing some
threshold. Scale changes det by $c^{n}$; it changes condition not at all.
\end{tripwire}

\section{Summary}

\begin{center}
\begin{tabularx}{\textwidth}{l>{\raggedright\arraybackslash}X}
\toprule
\textbf{Object} & \textbf{What you need to know}\\
\midrule
Cofactor expansion & $\det A=\sum_j a_{1j}(-1)^{1+j}M_{1j}$; any row or
column works; $\sim n!$ cost \\
Pivot product & $\det A=(-1)^{s}\prod u_{jj}$; the practical algorithm
behind \texttt{np.linalg.det} \\
Signed volume & columns of $A$ span a parallelepiped of volume $|\det A|$;
sign = orientation \\
Product rule & $\det(AB)=\det A\det B$; $\det(A^{-1})=1/\det A$;
$\det(cA)=c^{n}\det A$ \\
Row operations & add: unchanged; swap: $\times(-1)$; scale by $c$:
$\times c$ \\
$\det A=0$ & collapse: dependent columns, singular, no inverse --- all one
fact \\
Small det & not singular: conditioning is the question, use
\texttt{np.linalg.cond} \\
Cramer's rule & $x_j=\det A_j/\det A$: a theorem that produces a formula,
not an algorithm \\
\bottomrule
\end{tabularx}
\end{center}

\section*{Additional exercises}
\addcontentsline{toc}{section}{Additional exercises}

\begin{exercise}[easy]
Compute $\det\begin{bmatrix}1&4&2\\0&3&1\\0&0&5\end{bmatrix}$ --- and state
in one word why it is fast.
\end{exercise}

\begin{exercise}[easy]
$\det A=3$. What are $\det(2A)$, $\det(A^{T})$, $\det(A^{-1})$ for
$A\in\mathbb R^{3\times3}$?
\end{exercise}

\begin{exercise}[easy]
A $2\times2$ map sends the unit square to a parallelogram of area $7$.
Give two matrices with $|\det|=7$, one preserving and one flipping
orientation.
\end{exercise}

\begin{exercise}[easy]
True or false, justify in one line: ``$\det A>1$ means $A$ expands every
vector.''
\end{exercise}

\begin{exercise}[medium]
Prove $\det(A^{T})=\det A$ for $2\times2$ and $3\times3$ by direct
expansion, and explain in one sentence why it holds in general (expansion
along a column mirrors expansion along a row).
\end{exercise}

\begin{exercise}[medium]
Without computing, order these determinants of $3\times3$ matrices: $A$;
$B$ = $A$ with rows 1 and 3 swapped; $C$ = $A$ with row 2 doubled;
$D$ = $A$ with row 3 replaced by row 3 $+ 5\times$ row 1. Write each as a
multiple of $\det A$.
\end{exercise}

\begin{exercise}[medium]
In code: take the $4\times4$ Hilbert matrix \texttt{scipy.linalg.hilbert(4)}
(or build $H_{ij}=1/(i+j-1)$). Compute \texttt{np.linalg.det},
\texttt{np.linalg.cond}, and our \texttt{pivot\_det}
(\texttt{code/ch06\_pivot\_det.py}). Is $H$ singular? Would you trust det
to answer that?
\end{exercise}

\begin{exercise}[medium]
Use Cramer's rule to solve
$\begin{bmatrix}2&1\\1&3\end{bmatrix}\mathbf x=\begin{bmatrix}4\\5\end{bmatrix}$,
then solve the same system by elimination and compare the work.
\end{exercise}

\begin{exercise}[hard]
Prove $\det(AB)=\det A\det B$ for invertible $A$, general $B$, by
elimination: write $A=LU$ (ch.~2), show $\det(LU B)$ splits into
$\det L\det U\det B$ factors, and handle $\det B=0$ as a limiting case.
\end{exercise}

\begin{exercise}[hard]
A ``det'' candidate $f:\mathbb R^{n\times n}\to\mathbb R$ is required to be
multilinear in the columns, alternating (zero on equal columns), and
normalized $f(I)=1$. For $2\times2$ matrices only, prove
$f(A)=ad-bc$: expand each column as $a_{1j}\mathbf e_1+a_{2j}\mathbf e_2$,
distribute, and show every term except the two standard permutations
dies. (This is Axler's definition-first route, in miniature.)
\end{exercise}

\chapternotes{%
The cofactors/pivots/volume trio follows Strang's treatment in \cite{ila}
ch.~7 and the 18.06SC Unit II lectures \cite{ocw1806sc}; the
signed-volume pedagogy is the \emph{Essence of Linear Algebra} episode on
determinants \cite{eola}. Axler \cite{axler} ch.~10 builds the determinant
definition-first via multilinear forms --- the opposite order from ours,
worth reading once the three definitions here feel solid (the last
[hard] exercise rehearses his approach in $2\times2$). MML ch.~4
\cite{mml} frames the same properties for ML use.}
